\section*{Bloque B. Aplicaciones cuantitativas y optimización}

\noindent\textit{Preguntas 31 a 65. Evalúan la capacidad de resolver analíticamente equilibrios en cantidades y precios (Cournot, Bertrand, Stackelberg), funciones de mejor respuesta, matrices de juegos 2x2, factores de descuento críticos para colusión, exceso de capacidad y diferenciación de productos.}

\q{Considere un duopolio de Cournot simétrico con demanda inversa $P = 100 - Q$ (donde $Q = q_1 + q_2$) y costos marginales constantes $c_1 = c_2 = 10$. ¿Cuál es la cantidad de equilibrio producida por cada empresa?}
\begin{opts}
    \item $q_1 = q_2 = 20$
    \item $q_1 = q_2 = 45$
    \item $q_1 = q_2 = 25$
    \item $q_1 = q_2 = 30$
\end{opts}

\q{Para el mismo mercado de la pregunta anterior ($P = 100 - Q$, $c=10$), determine el precio de mercado de equilibrio y los beneficios individuales de cada empresa:}
\begin{opts}
    \item $P = 40$ y $\pi_1 = \pi_2 = 900$
    \item $P = 55$ y $\pi_1 = \pi_2 = 1012.5$
    \item $P = 70$ y $\pi_1 = \pi_2 = 1250$
    \item $P = 30$ y $\pi_1 = \pi_2 = 600$
\end{opts}

\q{En un oligopolio de Cournot con $n$ empresas simétricas, demanda inversa $P = a - bQ$ y costo marginal constante $c$ ($a > c$), la producción agregada total de la industria $Q^*$ es:}
\begin{opts}
    \item $Q^* = \frac{n(a - c)}{2b(n-1)}$
    \item $Q^* = \frac{n(a - c)}{b(n+2)}$
    \item $Q^* = \frac{n(a - c)}{b(n+1)}$
    \item $Q^* = \frac{(n+1)(a - c)}{bn}$
\end{opts}

\q{A partir de la fórmula general de Cournot con $n$ empresas ($P = a - bQ$, costo $c$), si la demanda es $P = 150 - Q$, el costo es $c = 30$ y operan $n = 5$ empresas, ¿cuál es el precio de mercado resultante?}
\begin{opts}
    \item $P = 60$
    \item $P = 50$
    \item $P = 40$
    \item $P = 30$
\end{opts}

\q{Considere un duopolio de Cournot con costos asimétricos: demanda $P = 120 - (q_1 + q_2)$, costo marginal de la empresa 1 $c_1 = 20$ y de la empresa 2 $c_2 = 40$. Las cantidades de equilibrio $(q_1^*, q_2^*)$ son:}
\begin{opts}
    \item $q_1^* = 45, q_2^* = 15$
    \item $q_1^* = 50, q_2^* = 10$
    \item $q_1^* = 30, q_2^* = 30$
    \item $q_1^* = 40, q_2^* = 20$
\end{opts}

\q{En el mercado asimétrico anterior ($q_1^* = 40, q_2^* = 20, P = 60$), los beneficios de ambas empresas $(\pi_1^*, \pi_2^*)$ son respectivamente:}
\begin{opts}
    \item $\pi_1^* = 1600$ y $\pi_2^* = 400$
    \item $\pi_1^* = 2000$ y $\pi_2^* = 500$
    \item $\pi_1^* = 1200$ y $\pi_2^* = 600$
    \item $\pi_1^* = 1400$ y $\pi_2^* = 800$
\end{opts}

\q{Considere un duopolio de Stackelberg con demanda inversa $P = 100 - Q$ y costos marginales constantes $c_1 = c_2 = 10$. La empresa 1 es el líder y la empresa 2 es la seguidora. ¿Cuánto produce el líder ($q_L^*$)?}
\begin{opts}
    \item $q_L^* = 40$
    \item $q_L^* = 30$
    \item $q_L^* = 45$
    \item $q_L^* = 50$
\end{opts}

\q{Siguiendo el caso de Stackelberg anterior ($q_L^* = 45$), la cantidad producida por la seguidora ($q_S^*$) y el precio de mercado ($P$) son:}
\begin{opts}
    \item $q_S^* = 20.0$ y $P = 35.0$
    \item $q_S^* = 15.0$ y $P = 40.0$
    \item $q_S^* = 30.0$ y $P = 25.0$
    \item $q_S^* = 22.5$ y $P = 32.5$
\end{opts}

\q{En el equilibrio de Stackelberg anterior ($q_L = 45, q_S = 22.5, P = 32.5, c = 10$), los beneficios $(\pi_L, \pi_S)$ son:}
\begin{opts}
    \item $\pi_L = 900.00$ y $\pi_S = 900.00$
    \item $\pi_L = 1200.0$ y $\pi_S = 400.00$
    \item $\pi_L = 1012.5$ y $\pi_S = 506.25$
    \item $\pi_L = 1125.0$ y $\pi_S = 625.00$
\end{opts}

\q{En un duopolio de Bertrand con bienes homogéneos, demanda de mercado $P = 180 - Q$ y costos marginales constantes $c_1 = c_2 = 30$, el precio y los beneficios de equilibrio son:}
\begin{opts}
    \item $P = 105$ y $\pi_1 = \pi_2 = 2812.5$
    \item $P = 30$ y $\pi_1 = \pi_2 = 0$
    \item $P = 80$ y $\pi_1 = \pi_2 = 1250$
    \item $P = 0$ y $\pi_1 = \pi_2 = -450$
\end{opts}

\q{Considere un duopolio de Bertrand con costos asimétricos: $c_1 = 25$ y $c_2 = 35$, con demanda $P = 100 - Q$. Si las empresas eligen precios en una grilla discreta de centavos, el precio fijado por la empresa 1 y su beneficio aproximado son:}
\begin{opts}
    \item $P = 62.5$ y $\pi_1 \approx 1406$
    \item $P = 34.99$ y $\pi_1 \approx 650$
    \item $P = 35$ y $\pi_1 = 0$
    \item $P = 25$ y $\pi_1 = 0$
\end{opts}

\q{Dos empresas compiten en precios con productos diferenciados. Las demandas son $q_1 = 24 - 2p_1 + p_2$ y $q_2 = 24 - 2p_2 + p_1$. Si los costos marginales son nulos ($c=0$), la función de reacción en precios de la empresa 1 es:}
\begin{opts}
    \item $p_1 = 12 - 0.5 p_2$
    \item $p_1 = 24 - 2 p_2$
    \item $p_1 = 8 + 0.5 p_2$
    \item $p_1 = 6 + 0.25 p_2$
\end{opts}

\q{En el sistema diferenciado de la pregunta anterior ($p_1 = 6 + 0.25 p_2$, con simetría $p_1 = p_2$), ¿cuál es el precio de equilibrio y el beneficio de cada empresa?}
\begin{opts}
    \item $p^* = 12$ y $\pi_1^* = \pi_2^* = 144$
    \item $p^* = 6$ y $\pi_1^* = \pi_2^* = 72$
    \item $p^* = 10$ y $\pi_1^* = \pi_2^* = 150$
    \item $p^* = 8$ y $\pi_1^* = \pi_2^* = 128$
\end{opts}

\q{Considere una industria con demanda $P = 140 - Q$ y costo marginal constante $c = 20$. Si dos empresas idénticas forman un cártel colusorio perfecto y se dividen la producción en partes iguales, la producción de cada una ($q_i^M$) y su beneficio colusorio ($\pi_i^M$) son:}
\begin{opts}
    \item $q_i^M = 25$ y $\pi_i^M = 1250$
    \item $q_i^M = 30$ y $\pi_i^M = 1800$
    \item $q_i^M = 35$ y $\pi_i^M = 1225$
    \item $q_i^M = 40$ y $\pi_i^M = 1600$
\end{opts}

\q{En el cártel anterior ($q_1 = q_2 = 30, c = 20, P = 140 - Q$), si la empresa 2 respeta su cuota $q_2 = 30$, ¿cuál es la cantidad óptima de desvío de la empresa 1 ($q_1^{dev}$) y su beneficio por hacer trampa ($\pi_1^{dev}$)?}
\begin{opts}
    \item $q_1^{dev} = 50$ y $\pi_1^{dev} = 2500$
    \item $q_1^{dev} = 45$ y $\pi_1^{dev} = 2025$
    \item $q_1^{dev} = 40$ y $\pi_1^{dev} = 1600$
    \item $q_1^{dev} = 35$ y $\pi_1^{dev} = 1925$
\end{opts}

\q{En el mismo duopolio ($P = 140 - Q, c = 20$), el beneficio de Cournot es $\pi_C = 1600$. Si las empresas usan la estrategia del gatillo (grim trigger), ¿cuál es el factor de descuento crítico $\delta^*$ mínimo para sostener la colusión?}
\begin{opts}
    \item $\delta^* = 3/4 \approx 0.750$
    \item $\delta^* = 1/2 \approx 0.500$
    \item $\delta^* = 9/17 \approx 0.529$
    \item $\delta^* = 2/3 \approx 0.667$
\end{opts}

\q{En un oligopolio de Bertrand simétrico con $n = 4$ empresas homogéneas que coluden al precio monopólico y usan grim trigger (donde la desviación captura todo el mercado y el castigo de Bertrand da beneficio cero), el factor de descuento mínimo requerido $\delta^*$ es:}
\begin{opts}
    \item $\delta^* = 0.50$
    \item $\delta^* = 0.75$
    \item $\delta^* = 0.80$
    \item $\delta^* = 0.25$
\end{opts}

\q{Considere la matriz de pagos de un juego simultáneo 2x2 entre dos empresas:
Firma 1 / Firma 2: Precio Alto (PA), Precio Bajo (PB)
PA: (100, 100), (20, 140)
PB: (140, 20), (50, 50)
¿Cuál es el único Equilibrio de Nash en estrategias puras de este juego?}
\begin{opts}
    \item (PA, PA)
    \item (PB, PA)
    \item (PB, PB)
    \item (PA, PB)
\end{opts}

\q{En el juego de coordinación 'Batalla de los sexos' con la siguiente matriz:
J1 / J2: Opción A, Opción B
Opción A: (2, 1), (0, 0)
Opción B: (0, 0), (1, 2)
¿Con qué probabilidad $p$ debe jugar J1 la Opción A en el equilibrio de Nash en estrategias mixtas?}
\begin{opts}
    \item $p = 1/2$
    \item $p = 3/4$
    \item $p = 1/3$
    \item $p = 2/3$
\end{opts}

\q{Considere un juego secuencial de entrada. Un Entrante potencial decide primero si Entrar ($E$) o No Entrar ($NE$). Si elige $NE$, los pagos son $(0, 80)$. Si elige $E$, el Incumbente decide si Luchar ($L$) o Acomodar ($A$). Si Lucha, los pagos son $(-10, 20)$; si Acomoda, los pagos son $(15, 40)$. ¿Cuál es el Equilibrio de Nash Perfecto en Subjuegos (SPNE)?}
\begin{opts}
    \item (NE; Lucha si entra), con pagos (0, 80)
    \item (E; Lucha), con pagos (-10, 20)
    \item (NE; Acomoda), con pagos (0, 80)
    \item (E; Acomoda), con pagos (15, 40)
\end{opts}

\q{En un modelo de competencia monopolística a corto plazo, la curva de demanda percibida por una empresa es $q_i = 60 - 2p_i + 2\bar{p}$ y su costo total es $CT(q_i) = 10q_i$. Si en equilibrio simétrico todos cobran el mismo precio ($p_i = \bar{p}$), ¿cuál es el precio de equilibrio de corto plazo?}
\begin{opts}
    \item $p^* = 25$
    \item $p^* = 15$
    \item $p^* = 20$
    \item $p^* = 30$
\end{opts}

\q{En una industria bajo competencia monopolística, cada firma enfrenta la curva de demanda $q = \frac{S}{n} - b(p - \bar{p})$, con costo total $CT(q) = F + c q$. A largo plazo, el número de empresas de equilibrio $n^*$ satisface la condición de beneficio nulo ($P = CMe$). La expresión para $n^*$ es:}
\begin{opts}
    \item $n^* = \frac{S}{F + c}$
    \item $n^* = \frac{S \cdot b}{F}$
    \item $n^* = \sqrt{\frac{S}{b F}}$
    \item $n^* = \sqrt{\frac{F}{S \cdot b}}$
\end{opts}

\q{La función de costo total de una empresa en competencia monopolística es $CT(q) = 72 + 2q + 0.5 q^2$. La escala mínima eficiente (donde el costo medio es mínimo) y la cantidad de largo plazo si produce con exceso de capacidad al nivel $q = 8$ implican un exceso de capacidad de:}
\begin{opts}
    \item 4 unidades
    \item 6 unidades
    \item 8 unidades
    \item 2 unidades
\end{opts}

\q{En el modelo de Hotelling sobre un segmento de longitud 1, dos firmas se ubican en los extremos 0 y 1. Los costos de transporte para los consumidores son $t x^2$ (cuadráticos) y los costos marginales de producción son $c = 0$. El precio de equilibrio simétrico fijado por ambas empresas es:}
\begin{opts}
    \item $p^* = t$
    \item $p^* = 2t$
    \item $p^* = \sqrt{t}$
    \item $p^* = t/2$
\end{opts}

\q{Para el modelo de Hotelling de la pregunta anterior ($p_1^* = p_2^* = t$, costo $c=0$, masa total de consumidores igual a 1), el beneficio de equilibrio de cada firma es:}
\begin{opts}
    \item $\pi_i^* = 2t$
    \item $\pi_i^* = t/4$
    \item $\pi_i^* = t$
    \item $\pi_i^* = t/2$
\end{opts}

\q{En el modelo de ciudad circular de Salop con circunferencia unitaria, costo de transporte $t x$ y costo marginal $c$, si operan $n$ empresas equidistantes, el precio de equilibrio simétrico es:}
\begin{opts}
    \item $p^* = c + \frac{t}{2n}$
    \item $p^* = c + t \cdot n$
    \item $p^* = c + \frac{t}{n}$
    \item $p^* = \frac{c + t}{n}$
\end{opts}

\q{En el modelo de variaciones conjeturales de Bowley, si la empresa 1 conjetura que $\frac{dq_2}{dq_1} = -1$, su ingreso marginal percibido bajo demanda inversa lineal $P(Q) = a - b(q_1 + q_2)$ coincide con:}
\begin{opts}
    \item un nivel estrictamente negativo en toda la región factible.
    \item el precio de mercado $P$, actuando como tomadora de precios.
    \item el ingreso marginal convencional de un duopolio de Cournot.
    \item el ingreso marginal correspondiente a un monopolista puro.
\end{opts}

\q{Considere un cártel donde el castigo por desvío dura exactamente $T = 1$ periodo (estrategia de reversión finita a Bertrand) en lugar de infinito. El factor de descuento mínimo requerido para coludir en duopolio simétrico ($n=2$) frente al grim trigger:}
\begin{opts}
    \item es estrictamente mayor, dificultando la colusión.
    \item es menor, facilitando la disciplina del cártel.
    \item es algebraicamente idéntico al del castigo perpetuo.
    \item se reduce a cero cuando la demanda es lineal.
\end{opts}

\q{En un modelo de empresa dominante con franja competitiva, la demanda total es $D(p) = 200 - p$ y la oferta agregada de la franja competitiva es $S_f(p) = p$. La curva de demanda residual que enfrenta la empresa dominante $D_R(p)$ es:}
\begin{opts}
    \item $D_R(p) = 200 - 2p$
    \item $D_R(p) = 100 - p$
    \item $D_R(p) = 200$
    \item $D_R(p) = 100 - 0.5p$
\end{opts}

\q{Siguiendo el caso de la empresa dominante anterior ($D_R(p) = 200 - 2p$, equivalente a $p = 100 - 0.5 q_L$), si el costo marginal de la empresa dominante es $c_L = 20$, ¿qué precio fijará en el mercado?}
\begin{opts}
    \item $P = 50$
    \item $P = 70$
    \item $P = 40$
    \item $P = 60$
\end{opts}

\q{La relación entre el índice de Lerner ponderado por cuotas de mercado ($L$) en un oligopolio de Cournot, el índice de Herfindahl-Hirschman ($HHI$) y la elasticidad precio de la demanda ($\varepsilon$) está dada por:}
\begin{opts}
    \item $L = HHI \cdot \varepsilon$
    \item $L = \frac{HHI}{\varepsilon}$
    \item $L = \sqrt{HHI} + \varepsilon$
    \item $L = \frac{\varepsilon}{HHI}$
\end{opts}

\q{En una industria con demanda lineal y costos marginales constantes, la pérdida irrecuperable de eficiencia ($DWL$) en un duopolio simétrico de Cournot comparada con la del monopolio puro ($DWL_M$) equivale exactamente a:}
\begin{opts}
    \item $\frac{2}{3} DWL_M$
    \item $\frac{1}{4} DWL_M$
    \item $\frac{4}{9} DWL_M$
    \item $\frac{1}{2} DWL_M$
\end{opts}

\q{Para la demanda $P = 100 - Q$ y costo $c = 10$, el Excedente del Consumidor ($EC$) en el duopolio de Cournot ($Q = 60$) frente al monopolio ($Q = 45$) es:}
\begin{opts}
    \item $EC_{Cournot} = 1600.0$ vs $EC_{Monopolio} = 900.00$
    \item $EC_{Cournot} = 2000.0$ vs $EC_{Monopolio} = 1250.0$
    \item $EC_{Cournot} = 1800.0$ vs $EC_{Monopolio} = 1012.5$
    \item $EC_{Cournot} = 2400.0$ vs $EC_{Monopolio} = 1500.0$
\end{opts}

\q{Si se introduce un impuesto específico de $t$ soles por unidad en un duopolio de Cournot simétrico con demanda lineal $P = a - b Q$, la traslación del impuesto al precio final del consumidor ($dP/dt$) es:}
\begin{opts}
    \item $dP/dt = 1$
    \item $dP/dt = 1/3$
    \item $dP/dt = 1/2$
    \item $dP/dt = 2/3$
\end{opts}

\q{La 'Paradoja de las Fusiones' demostrada por Salant, Switzer y Reynolds (1983) en un oligopolio de Cournot simétrico con demanda lineal y sin sinergias de costos establece que:}
\begin{opts}
    \item una fusión horizontal bilateral entre competidores simétricos siempre aumenta los beneficios conjuntos.
    \item una fusión bilateral reduce los beneficios conjuntos frente a sus rivales, salvo fusiones casi totales.
    \item las fusiones horizontales elevan siempre el excedente del consumidor aun sin sinergias de eficiencia.
    \item las firmas competidoras no fusionadas responden contrayendo drásticamente su nivel de producción.
\end{opts}

